% ============================================================
%  Section 4 — Results
% ============================================================
\section*{Results}

\subsection*{4.1 Experimental Overview}

Eight model configurations were evaluated on a held-out test subject (case~000) not seen during
training.  The configurations span the full factorial combination of two task modes, two network
architectures, and two input conditioning strategies.  In \textbf{Denoising} mode ($\times 1$),
the network translates clinical 3T acquisitions into the contrast and noise characteristics of
7T, operating at native resolution.  In \textbf{Super-Resolution} mode ($\times 2$, hereafter
\textbf{SR}), the network simultaneously denoises and upsamples by a factor of two, reconstructing
7T-equivalent volumes from $\times 0.5$-downsampled 3T inputs.  The two architectures are the
\textbf{Original} SR4DFlowNet baseline and the \textbf{CBAM-PreUpsampling} extension, which
inserts a convolutional block attention module before the final upsampling stage.  The two input
strategies are \textbf{Masked} training, in which ground-truth supervision is restricted to the
Circle of Willis mask $\mathcal{M}$, and \textbf{Synthetic-noise} training, in which an unmasked
3T input augmented with simulated acquisition noise is used.

Table~\ref{tab:experiment_map} maps each identifier to its configuration.  All geometry metrics
are computed against a common 7T reference segmentation of 15{,}577 voxels at 0.5\,mm isotropic
resolution.  Velocity and background metrics are averaged across nine cardiac frames, with all
velocity values expressed in VENC-normalised units.

\begin{table}[h!]
  \centering
  \caption{Summary of the eight evaluated configurations.}
  \label{tab:experiment_map}
  \small
  \begin{tabular}{clll}
    \toprule
    ID & Task mode & Architecture & Input strategy \\
    \midrule
    A & Denoising ($\times 1$) & Original           & Masked \\
    B & Denoising ($\times 1$) & Original           & Synth.\ noise \\
    C & Denoising ($\times 1$) & CBAM-PreUpsampling & Masked \\
    D & Denoising ($\times 1$) & CBAM-PreUpsampling & Synth.\ noise \\
    \midrule
    E & Super-Resolution ($\times 2$) & Original           & Masked \\
    F & Super-Resolution ($\times 2$) & Original           & Synth.\ noise \\
    G & Super-Resolution ($\times 2$) & CBAM-PreUpsampling & Masked \\
    H & Super-Resolution ($\times 2$) & CBAM-PreUpsampling & Synth.\ noise \\
    \bottomrule
  \end{tabular}
\end{table}

Performance is assessed along three complementary axes: (i)~\textbf{velocity reconstruction
fidelity} via the masked MSE of all three velocity components jointly ($\text{MSE}_\text{vel}$)
and per-directional breakdown ($u$: feet--head, $v$: left--right, $w$: anterior--posterior);
(ii)~\textbf{vascular geometry fidelity} via the Dice coefficient, the symmetric mean surface
distance (sMSD), and the Hausdorff distance (HD) of the re-segmented prediction mask against the
7T reference; and (iii)~\textbf{background suppression} via the MSE of predicted velocities
outside $\mathcal{M}$ ($\text{MSE}_\text{bg}$).

% -------------------------------------------------------
\subsection*{4.2 Velocity Reconstruction}
% -------------------------------------------------------

Table~\ref{tab:velocity} reports velocity MSE within the CoW mask.  Bold denotes the best value
in each column within each task mode.

\begin{table}[h!]
  \centering
  \caption{Intraluminal velocity reconstruction metrics.  All values are VENC-normalised MSE
  averaged over nine cardiac frames (lower is better).}
  \label{tab:velocity}
  \setlength{\tabcolsep}{5pt}
  \small
  \begin{tabular}{cll l rrrr}
    \toprule
    \multirow{2}{*}{ID} & \multirow{2}{*}{Architecture} & \multirow{2}{*}{Input} &
      & \multicolumn{4}{c}{MSE $\downarrow$ (VENC-normalised)} \\
    \cmidrule(lr){5-8}
    & & & & $\text{Total}$ & $u$ & $v$ & $w$ \\
    \midrule
    \multicolumn{8}{l}{\textit{Denoising ($\times 1$)}} \\
    \midrule
    A & Original           & Masked       & & \textbf{0.04281} & 0.01354 & 0.01748 & 0.01180 \\
    B & Original           & Synth.\ noise & & 0.04721 & 0.01440 & 0.02054 & 0.01227 \\
    C & CBAM-PreUpsampling & Masked       & & 0.04306 & \textbf{0.01327} & 0.01815 & \textbf{0.01164} \\
    D & CBAM-PreUpsampling & Synth.\ noise & & 0.04399 & 0.01458 & \textbf{0.01725} & 0.01216 \\
    \midrule
    \multicolumn{8}{l}{\textit{Super-Resolution ($\times 2$)}} \\
    \midrule
    E & Original           & Masked       & & 0.04977 & \textbf{0.01472} & 0.01853 & 0.01652 \\
    F & Original           & Synth.\ noise & & 0.04983 & 0.01592 & 0.01903 & 0.01489 \\
    G & CBAM-PreUpsampling & Masked       & & 0.04844 & \textbf{0.01472} & 0.01833 & 0.01539 \\
    H & CBAM-PreUpsampling & Synth.\ noise & & \textbf{0.04833} & 0.01577 & \textbf{0.01807} & \textbf{0.01449} \\
    \bottomrule
  \end{tabular}
\end{table}

In the Denoising mode, masked training retains a clear velocity advantage.  Configuration~A
achieves the lowest total MSE (0.04281), a 9.3\% reduction over its synthetic-noise counterpart~B
(0.04721).  The CBAM-PreUpsampling architecture narrows but does not close this gap:
configuration~D (0.04399) outperforms~B by 6.8\% and approaches the masked configurations, with
a residual deficit of only 2.7\% relative to~C (0.04306).  These results indicate that the
attention mechanism provides a substantial architectural benefit when mask-guided supervision is
unavailable.  The advantage of masked training under CBAM is accordingly smaller (2.1\%), as the
architecture already recovers much of the information that masking contributes.  At the
per-component level, configurations~C and~A trade leadership: CBAM-masked achieves the lowest
$u$- and $w$-direction MSE (0.01327 and 0.01164 respectively), while the original masked model~A
records the lowest total and $v$-direction MSE (0.01748).  Interestingly, configuration~D records
the best $v$-component of all Denoising models (0.01725), slightly outperforming~A, suggesting
that for the left--right direction the CBAM attention mechanism provides a specific benefit even
without mask supervision.

In the Super-Resolution mode, a qualitatively different pattern emerges.  The input strategy has
almost no effect on total velocity MSE when the CBAM architecture is used: configuration~G
(masked) and~H (synth.\ noise) differ by only 0.00011 (0.2\%), with~H (0.04833) marginally
lower.  This near-zero gap implies that, for the harder $\times 2$ task, the pre-upsampling
attention module is able to recover the velocity accuracy that masked supervision would otherwise
provide.  For the Original architecture the gap is similarly minimal (E: 0.04977 vs.\ F:
0.04983).  The dominant factor in SR velocity accuracy is therefore the architecture rather than
the input strategy: CBAM reduces total MSE by 2.7\% over Original under masked training (G vs.\ E)
and by 3.0\% under synthetic noise (H vs.\ F).  Configuration~H also achieves the lowest
$v$- and $w$-component MSE across all SR models (0.01807 and 0.01449 respectively), reinforcing
that the CBAM architecture provides the principal source of velocity accuracy in this mode.

% -------------------------------------------------------
\subsection*{4.3 Vascular Geometry Fidelity}
% -------------------------------------------------------

Geometric fidelity metrics are reported in Table~\ref{tab:geometry}.  Each predicted 4D Flow
volume was re-segmented using the CoW segmentation pipeline applied to the original 7T data,
and the resulting mask was compared against the 7T reference (15{,}577 voxels).

\begin{table}[h!]
  \centering
  \caption{Vascular geometry metrics.  The re-segmented prediction mask is compared against the
  7T CoW reference (15{,}577 voxels at 0.5\,mm isotropic).  Dice: volumetric overlap
  ($\uparrow$); sMSD: symmetric mean surface distance in mm ($\downarrow$); HD: Hausdorff
  distance in mm ($\downarrow$).}
  \label{tab:geometry}
  \setlength{\tabcolsep}{4.5pt}
  \small
  \begin{tabular}{cll l r rrr}
    \toprule
    ID & Architecture & Input & & Pred.\ voxels &
      Dice $\uparrow$ & sMSD $\downarrow$ & HD $\downarrow$ \\
    \midrule
    \multicolumn{8}{l}{\textit{Denoising ($\times 1$) — reference: 15,577 voxels}} \\
    \midrule
    A & Original           & Masked        & & 18{,}027 & 0.843 & 0.202 & \textbf{8.57} \\
    B & Original           & Synth.\ noise & & 14{,}840 & 0.734 & 0.667 & 15.09 \\
    C & CBAM-PreUpsampling & Masked        & & 17{,}544 & \textbf{0.855} & \textbf{0.183} & 9.07 \\
    D & CBAM-PreUpsampling & Synth.\ noise & & 15{,}533 & 0.761 & 0.501 & 9.86 \\
    \midrule
    \multicolumn{8}{l}{\textit{Super-Resolution ($\times 2$) — reference: 15,577 voxels}} \\
    \midrule
    E & Original           & Masked        & & 18{,}623 & \textbf{0.715} & \textbf{0.412} & \textbf{9.01} \\
    F & Original           & Synth.\ noise & & 12{,}808 & 0.694 & 0.828 & 14.82 \\
    G & CBAM-PreUpsampling & Masked        & & 19{,}294 & 0.702 & 0.413 & 9.50 \\
    H & CBAM-PreUpsampling & Synth.\ noise & & 12{,}841 & 0.698 & 0.737 & 13.59 \\
    \bottomrule
  \end{tabular}
\end{table}

Geometry is the dimension most sensitive to training strategy.  In the Denoising mode, masked
configurations (A and~C) substantially outperform synthetic-noise counterparts on all three
metrics.  Configuration~C achieves the best Dice (0.855) and sMSD (0.183\,mm) of all Denoising
models.  Configuration~A records the best Hausdorff distance (8.57\,mm), 0.50\,mm lower than~C,
reflecting tighter containment of worst-case boundary outliers.  Both masked models predict
volumes closely bracketing the reference (17{,}544--18{,}027 vs.\ 15{,}577 voxels), indicative
of a modest, controlled over-segmentation of peripheral branches.

The benefit of the CBAM architecture in the synthetic-noise Denoising setting (D vs.\ B) is
striking: Dice improves from 0.734 to 0.761 (+3.7\,pp) and Hausdorff distance falls from
15.09 to 9.86\,mm, a 34.6\% reduction.  This confirms that the attention mechanism provides an
inductive bias that partially compensates for the absence of mask-guided supervision, keeping
worst-case boundary errors close to the masked-training range (HD range for masked:
8.57--9.07\,mm).

In the Super-Resolution mode, Dice scores are uniformly lower than in Denoising ($\approx$14\,pp
gap for the best configurations: E 0.715 vs.\ C 0.855), consistent with the harder task.  The
relative ordering by training strategy is preserved: masked configurations (E and~G) achieve
lower HD than their synthetic-noise counterparts (F and~H), with configurations~E and~G both
producing HD values below 10\,mm.  Among SR models, the Original with masked input (E) achieves
the best Dice (0.715), the best sMSD (0.412\,mm), and the lowest HD (9.01\,mm), marginally
outperforming CBAM-masked (G: Dice 0.702, sMSD 0.413\,mm, HD 9.50\,mm).  The larger predicted
voxel count of~G (19{,}294 vs.\ 18{,}623) points to slight over-segmentation as the explanation
for this Dice gap.  For synthetic-noise SR, CBAM (H) reduces the HD relative to the Original (F)
from 14.82 to 13.59\,mm ($-$8.3\%), a smaller but consistent pattern of architectural benefit.

% -------------------------------------------------------
\subsection*{4.4 Background Suppression}
% -------------------------------------------------------

Table~\ref{tab:background} reports background velocity MSE (outside $\mathcal{M}$) for all eight
configurations.

\begin{table}[h!]
  \centering
  \caption{Background velocity MSE outside the Circle of Willis mask.  Values are VENC-normalised
  mean squared errors averaged over nine cardiac frames (lower is better).}
  \label{tab:background}
  \setlength{\tabcolsep}{4.5pt}
  \small
  \begin{tabular}{cll l rrrr}
    \toprule
    ID & Architecture & Input & &
      $\text{MSE}_{\text{bg}}$ $\downarrow$ &
      $\text{MSE}_{u}$ $\downarrow$ &
      $\text{MSE}_{v}$ $\downarrow$ &
      $\text{MSE}_{w}$ $\downarrow$ \\
    \midrule
    \multicolumn{8}{l}{\textit{Denoising ($\times 1$)}} \\
    \midrule
    A & Original           & Masked       & & $6.44{\times}10^{-4}$ & $1.85{\times}10^{-4}$ & $3.23{\times}10^{-4}$ & $1.36{\times}10^{-4}$ \\
    B & Original           & Synth.\ noise & & $1.51{\times}10^{-3}$ & $3.57{\times}10^{-4}$ & $8.41{\times}10^{-4}$ & $3.13{\times}10^{-4}$ \\
    C & CBAM-PreUpsampling & Masked       & & $\mathbf{6.08{\times}10^{-4}}$ & $\mathbf{1.52{\times}10^{-4}}$ & $3.37{\times}10^{-4}$ & $\mathbf{1.19{\times}10^{-4}}$ \\
    D & CBAM-PreUpsampling & Synth.\ noise & & $1.61{\times}10^{-3}$ & $4.59{\times}10^{-4}$ & $8.55{\times}10^{-4}$ & $3.00{\times}10^{-4}$ \\
    \midrule
    \multicolumn{8}{l}{\textit{Super-Resolution ($\times 2$)}} \\
    \midrule
    E & Original           & Masked       & & $7.83{\times}10^{-4}$ & $1.98{\times}10^{-4}$ & $4.14{\times}10^{-4}$ & $1.71{\times}10^{-4}$ \\
    F & Original           & Synth.\ noise & & $1.50{\times}10^{-3}$ & $4.72{\times}10^{-4}$ & $6.91{\times}10^{-4}$ & $3.38{\times}10^{-4}$ \\
    G & CBAM-PreUpsampling & Masked       & & $\mathbf{6.71{\times}10^{-4}}$ & $\mathbf{1.74{\times}10^{-4}}$ & $\mathbf{3.31{\times}10^{-4}}$ & $\mathbf{1.66{\times}10^{-4}}$ \\
    H & CBAM-PreUpsampling & Synth.\ noise & & $1.42{\times}10^{-3}$ & $4.02{\times}10^{-4}$ & $7.53{\times}10^{-4}$ & $2.68{\times}10^{-4}$ \\
    \bottomrule
  \end{tabular}
\end{table}

Background suppression is almost entirely determined by the training strategy.  Masked
configurations reduce extravascular MSE by approximately 2.0--2.7$\times$ relative to their
synthetic-noise counterparts: A vs.\ B ($\times 2.35$), C vs.\ D ($\times 2.65$), E vs.\ F
($\times 1.92$), and G vs.\ H ($\times 2.12$).  This reflects the explicit mask-aware loss
formulation, which penalises background predictions through the total variation term
$\mathcal{L}_\text{TV}$; the synthetic-noise strategy provides no equivalent regularisation
outside $\mathcal{M}$.

Within masked configurations, CBAM-PreUpsampling achieves lower background MSE than the Original
in both modes (C vs.\ A; G vs.\ E), suggesting that the attention mechanism improves spatial
localisation of predicted flow and thereby reduces extravascular signal leakage.  Within
synthetic-noise configurations, the architecture-background relationship differs by mode: in
Denoising, configuration~D (CBAM) produces the highest background MSE of any model
($1.61{\times}10^{-3}$), exceeding its Original counterpart~B ($1.51{\times}10^{-3}$); in SR,
the ordering reverses and CBAM~H ($1.42{\times}10^{-3}$) is slightly better than Original~F
($1.50{\times}10^{-3}$).  This reversal is consistent with the observation from
Section~\ref{sec:velocity} that the CBAM architecture has a qualitatively stronger impact on
SR than on Denoising in the absence of mask supervision.

The $v$-component (left--right phase-contrast direction) is consistently the largest contributor
to background error across all eight configurations, reaching values approximately two to three
times higher than $u$ or $w$ in the synthetic-noise setting.  This systematic pattern is
consistent with the velocity reconstruction results and likely reflects residual phase errors
or partial-volume effects in the left--right encoding direction of the 3T acquisition.

% -------------------------------------------------------
\subsection*{4.5 Consolidated Comparison}
% -------------------------------------------------------

Table~\ref{tab:summary} consolidates the five primary performance indicators across all eight
configurations for direct cross-experiment comparison.

\begin{table}[h!]
  \centering
  \caption{Consolidated results.  Bold denotes the best value within each task mode for each
  metric.  $\Delta_\text{DEN}$: gap to the best Denoising result;
  $\Delta_\text{SR}$: gap to the best SR result (not shown separately but implied by bold).}
  \label{tab:summary}
  \setlength{\tabcolsep}{4pt}
  \small
  \begin{tabular}{cll l rrrrr}
    \toprule
    ID & Architecture & Input & &
      $\text{MSE}_\text{vel}$ $\downarrow$ &
      Dice $\uparrow$ &
      sMSD $\downarrow$ &
      HD $\downarrow$ &
      $\text{MSE}_\text{bg}$ $\downarrow$ \\
    \midrule
    \multicolumn{9}{l}{\textit{Denoising ($\times 1$)}} \\
    \midrule
    A & Original      & Masked       & & \textbf{0.04281} & 0.843 & 0.202 & \textbf{8.57} & $6.44{\times}10^{-4}$ \\
    B & Original      & Synth.\ noise & & 0.04721 & 0.734 & 0.667 & 15.09 & $1.51{\times}10^{-3}$ \\
    C & CBAM-PreUps.  & Masked       & & 0.04306 & \textbf{0.855} & \textbf{0.183} & 9.07 & $\mathbf{6.08{\times}10^{-4}}$ \\
    D & CBAM-PreUps.  & Synth.\ noise & & 0.04399 & 0.761 & 0.501 & 9.86 & $1.61{\times}10^{-3}$ \\
    \midrule
    \multicolumn{9}{l}{\textit{Super-Resolution ($\times 2$)}} \\
    \midrule
    E & Original      & Masked       & & 0.04977 & \textbf{0.715} & \textbf{0.412} & \textbf{9.01} & $7.83{\times}10^{-4}$ \\
    F & Original      & Synth.\ noise & & 0.04983 & 0.694 & 0.828 & 14.82 & $1.50{\times}10^{-3}$ \\
    G & CBAM-PreUps.  & Masked       & & 0.04844 & 0.702 & 0.413 & 9.50 & $\mathbf{6.71{\times}10^{-4}}$ \\
    H & CBAM-PreUps.  & Synth.\ noise & & \textbf{0.04833} & 0.698 & 0.737 & 13.59 & $1.42{\times}10^{-3}$ \\
    \bottomrule
  \end{tabular}
\end{table}

Three principal conclusions emerge from the consolidated results.

\textbf{Mask-guided training is the dominant factor for geometric fidelity and background
suppression.}  Across both task modes, masked configurations consistently achieve lower HD
(by 5--6\,mm relative to synthetic-noise counterparts using the same architecture) and lower
background MSE (by 1.9--2.7$\times$), while Dice scores improve by 9--11\,pp in Denoising and
up to 2\,pp in SR.  These gains are direct consequences of the mask-restricted loss and the
$\mathcal{L}_\text{TV}$ background penalty.  Notably, the Hausdorff distance is the metric most
sensitive to the training strategy: Original-architecture synthetic-noise training produces worst-case boundary errors above
14\,mm (B: 15.09, F: 14.82), while CBAM synthetic-noise training reduces HD to 13.59\,mm (H)
and even to 9.86\,mm (D), approaching the masked-training range; all masked models remain below
10\,mm regardless of architecture.

\textbf{The CBAM-PreUpsampling architecture is most impactful in the absence of mask supervision
and in the SR setting.}  In the masked Denoising setting, CBAM (C) improves Dice by 1.2\,pp and
sMSD by 9.4\% over the Original (A), while the velocity MSE change is negligible ($+$0.6\%).  The
CBAM benefit is substantially larger in the synthetic-noise Denoising setting: Dice improves from
0.734 to 0.761 (+3.7\,pp) and HD falls from 15.09 to 9.86\,mm ($-$34.6\%), consistent with the
attention mechanism providing a spatial inductive bias that partially compensates for the absence
of explicit mask guidance.  In SR, the architecture effect on velocity is the dominant signal:
CBAM reduces $\text{MSE}_\text{vel}$ by 2.7--3.0\% in both input conditions, and G achieves the
lowest background MSE for SR ($6.71{\times}10^{-4}$).

\textbf{Input strategy and architecture interact differently across task modes.}  For Denoising,
masked training provides a clear, consistent velocity advantage (9.3\% for Original; 2.1\% for
CBAM).  For SR, this advantage effectively vanishes for the CBAM architecture: the velocity MSE
difference between G (masked) and H (synth.\ noise) is only 0.2\%, with~H (0.04833) marginally
outperforming~G (0.04844) on total and per-component $v$ and $w$ errors.  This suggests that
the pre-upsampling attention mechanism can substitute for mask-guided supervision in preserving
intraluminal velocity accuracy when the task complexity is high, at the cost of worse geometric
fidelity (HD: 13.59 vs.\ 9.50\,mm) and background suppression ($1.42$ vs.\
$0.67 {\times}10^{-3}$).  The implication for practical deployment is that the choice of
training strategy remains critical for applications requiring accurate vessel boundary recovery
or clean extravascular predictions, even when velocity reconstruction accuracy alone might
suggest that the simpler synthetic-noise pipeline is sufficient.

\label{sec:velocity}

\bibliographystyle{ieeetr}
\bibliography{neuroflow}
